\section{Chaîne exceptionnelle, théorie conforme, cordes et classification des masses}
\Needspace{7\baselineskip}
\subsection{Antécédent commun et coordination des \(2\) lectures de la chaîne exceptionnelle}
La relation entre masses et symétries exceptionnelles ne commence pas par une action du Monstre sur une table de particules. Elle commence dans le même état enrichi du TPK. Soit

\[
x\in\mathcal X_{\rm TPK}^{\rm enr}.
\]
La lecture archimédienne contracte des cylindres compatibles et conserve la valeur réelle avec sa généalogie :

\[
\mathcal P_{\rm arq}(x)\in\mathbb R.
\]
La lecture exceptionnelle conserve les incidences de frontière et produit une classe dans la limite projective :

\[
\mathcal P_{\rm exc}(x)=Q_\infty(x).
\]
La naturalité exige, pour tout raffinement \(n\ge m\),

\[
r_{n,m}\circ Q_n=Q_m\circ p_{n,m}.
\]
L’application conjointe

\[
\mathcal P_{\rm joint}:x\longmapsto
\bigl(\mathcal P_{\rm arq}(x),\mathcal P_{\rm exc}(x)\bigr)
\]
est la forme positive de la double projection. Elle ne nie pas la connexion Monstre--masse : elle la type. La branche de masse et la branche exceptionnelle sont les images d’un même antécédent doté d’une orientation, d’une graduation, d’une retenue, d’une mémoire et d’une holonomie communes. Elle n’affirme pas non plus que le Monstre agit littéralement sur les chiffres décimaux ou qu’un coefficient de Moonshine est à lui seul une masse.

\Needspace{7\baselineskip}
\subsection{Construction de la chaîne exceptionnelle}
L’incidence ternaire prolonge l’état par la chaîne

\[
C_3
\longrightarrow A_2^{12}
\longrightarrow N(A_2^{12})
\longrightarrow \Lambda_{24}
\longrightarrow V^\natural
\longrightarrow \mathbb M,
\]
où \(\Lambda_{24}\) est le réseau de Leech, \(V^\natural\) l’algèbre de vertex du Moonshine et \(\mathbb M\) le groupe Monstre. Le passage ne consiste pas à juxtaposer des noms classiques : chaque flèche doit préserver l’incidence, la graduation et la compatibilité de raffinement héritées de l’état.

La trace ternaire s’écrit

\[
t_3(q)=q^{-1}
\prod_{n\ge1}(1+q^n+q^{2n})^{-12}.
\]
Le premier défaut non trivial vérifie

\[
D_{\rm APP}=[q]t_3=v_\alpha^2=54,
\]
et le premier coefficient du caractère du Monstre conserve la décomposition

\[
196884=54+10\cdot3^9.
\]
Ces identités transportent la graduation et la multiplicité. Elles n’attribuent pas de masse en MeV. Leur fonction dans la théorie des masses est de classifier l’espace sur lequel agit l’opérateur avant de le scalariser.

\Needspace{7\baselineskip}
\subsection{Appariement fibre à fibre avec la loi des masses}
La branche matérielle du même état produit

\[
x\longmapsto
\bigl(F_X,\sigma_X,\eta_X,\widehat M_X\bigr),
\]
où \(F_X\) est un point, une orbite, une fibre ou une multisection ; \(\sigma_X\in\mathbb Z^5\) est la signature ; \(\eta_X\in\mathcal E_{90,120}\) est le raffinement ; et \(\widehat M_X\) l’opérateur de masse. La lecture conjointe est alors

\[
\mathfrak J_{\rm mass,exc}(x)
=\bigl(F_X,\sigma_X,\eta_X,\widehat M_X,Q_\infty(x)\bigr).
\]
Deux états ayant la même valeur scalaire de masse peuvent rester distincts par leur composante exceptionnelle ; deux états ayant le même caractère exceptionnel peuvent posséder des opérateurs de masse différents. La classification correcte est l’appariement des fibres, et non l’égalité isolée des nombres.

Si une représentation exceptionnelle \(G\) agit dans une fibre physique \(\mathcal H_X\), la compatibilité avec la masse exige

\[
U_g\widehat M_XU_g^{-1}=\widehat M_X,
\qquad g\in G.
\]
La décomposition

\[
\mathcal H_X=\bigoplus_{\lambda}
\mathcal H_{X,\lambda}
\]
sépare les multiplets exceptionnels. Sur un facteur direct irréductible, la commutation impose que la compression de \(\widehat M_X\) soit scalaire. C’est une voie légitime de scalarisation par symétrie. La valeur du scalaire continue de provenir du caractère, des tours, de la section dimensionnelle et du couplage ; la dimension de la représentation ne détermine que la dégénérescence.

\Needspace{7\baselineskip}
\subsection{Fock et algèbre de vertex}
Le prolongement d’un parcours à un espace à plusieurs particules possède deux complétions principales. La complétion extérieure produit les CAR, Pauli et la statistique fermionique ; la complétion symétrique produit les CCR et la statistique bosonique. Une algèbre de vertex ajoute un produit d’opérateurs, un vide, une translation et une graduation :

\[
V=\bigoplus_{n\ge0}V_n,
\qquad
L_0|_{V_n}=nI.
\]
Le caractère

\[
\operatorname{ch}_V(q)
=\operatorname{Tr}_V q^{L_0-c/24}
\]
compte les états par degré. Une multiplicité dans \(V_n\) n’est pas une masse. Pour transformer la graduation en énergie, il faut un Hamiltonien et une section de longueur ou de temps. Dans une réalisation conforme sur un cylindre de circonférence \(L\),

\[
H_{\rm CFT}
=\frac{2\pi\hbar_{\rm int}v}{L}
\left(L_0+\bar L_0-\frac{c}{12}\right).
\]
La masse ou le gap physique n’apparaît qu’après fixation de \(v\), de \(L\), des conditions de bord et de l’application du parcours TPK vers les modules conformes. La fonction de la chaîne exceptionnelle consiste à organiser les dégénérescences et les règles de composition de ces modules.

\Needspace{7\baselineskip}
\subsection{Spectre de corde comme réalisation dimensionnelle ultérieure}
Soit \(\alpha'\in\mathcal L_L^2\) la longueur quadratique représentant la tension inverse dans la carte naturalisée de corde déclarée. Dans une compactification circulaire de rayon \(R\), les impulsions gauche et droite sont

\[
p_L=\frac{m}{R}+\frac{wR}{\alpha'},
\qquad
p_R=\frac{m}{R}-\frac{wR}{\alpha'},
\]
avec \(m,w\in\mathbb Z\). Dans les unités naturelles de la réalisation, le spectre fermé vérifie

\[
M^2=p_L^2+\frac{4}{\alpha'}(N_L-a_L)
=p_R^2+\frac{4}{\alpha'}(N_R-a_R),
\]
ainsi que la condition de niveau

\[
N_L-N_R=a_L-a_R-mw.
\]
L'involution

\[
(m,w,R)\longleftrightarrow
\left(w,m,\frac{\alpha'}{R}\right)
\]
conserve le spectre : c’est la réalisation de la dualité T. HMT apporte au préalable l’orientation bidirectionnelle, la mémoire et l’appariement des feuillets ; la théorie des cordes ajoute la tension, le rayon, les oscillateurs, le niveau et les conditions aux limites. On n’obtient pas une masse de corde en remplaçant \(N_L,N_R\) par des hauteurs de la tour \(\langle90,120\rangle\). Toute identification entre les deux graduations exige un opérateur d’entrelacement qui conserve la somme, l’orientation et le caractère.

\Needspace{7\baselineskip}
\subsection{Écrans réticulaires de théorie M : graduation, compactification et transport de la mémoire résiduelle}
Les écrans internes

\[
12\longrightarrow11\longrightarrow10,
\qquad
26\longrightarrow24,
\]
conservent la descente dodécaphasique, l’élimination d’une direction redondante et l’écran transversal. Dans une réalisation de théorie M, les champs

\[
g_{MN},\qquad C_{MNP},\qquad\Psi_M
\]
doivent être obtenus au moyen d’un opérateur d’entrelacement préservant la jauge, la supersymétrie, les contraintes et le produit scalaire. Les écrans numériques ne remplacent pas cette application.

La classification des masses est transportée vers cette branche au moyen de

\[
\text{route}
\longrightarrow\text{registre généalogique}
\longrightarrow\text{signature}
\longrightarrow\text{caractère}
\longrightarrow\text{opérateur}
\longrightarrow\text{spectre compactifié}.
\]
Les modes de Kaluza--Klein, d’enroulement, de membrane ou d’excitation oscillatoire sont des états distincts même s’ils partagent une masse. Le graviton d’une réalisation de corde fermée, s’il apparaît, est un mode collectif sans masse du secteur géométrique ; il ne devient pas pour autant une graine élémentaire du TPK et n’invalide pas la formulation de la gravité en termes d’holonomie et de torsion.

\Needspace{7\baselineskip}
\subsection{Réalisations ordinaires et hypothétiques induites par la chaîne exceptionnelle}
La branche exceptionnelle permet de distinguer quatre sources de multiplicité :

\begin{enumerate}
\item multiplicité cellulaire ou de parcours dans l’atlas ;
\item dégénérescence d’une représentation exceptionnelle ;
\item occupation de Fock ;
\item niveau conforme, oscillatoire ou compactifié.
\end{enumerate}
Les confondre produit de faux partenaires, de fausses générations ou des masses dupliquées. Un état jumeau, supersymétrique, sombre ou exotique ne constitue une nouvelle espèce que lorsque son quintuplet

\[
(F_X,Q_\infty(x),\rho_X,\widehat M_X,P_X)
\]
diffère physiquement et que l’Hamiltonien possède un pôle ou une section persistante supplémentaire. Une coïncidence de dimension, de caractère ou de masse ne suffit pas.

Inversement, une dégénérescence protégée par la chaîne exceptionnelle est une prédiction structurelle : les perturbations préservant la symétrie ne peuvent scinder le multiplet. Une scission observée exige l’identification de la brisure de symétrie, du couplage ou de la carte qui la produit.

\Needspace{7\baselineskip}
\subsection{Conditions de positivité, de clôture et de stabilité des réalisations topologiques}
La connexion entre la chaîne exceptionnelle et les masses est réfutée dans une réalisation concrète si l’un des faits suivants se produit :

\begin{itemize}
\item la naturalité \(r_{n,m}Q_n=Q_mp_{n,m}\) échoue ;
\item l’étiquette exceptionnelle change lorsqu’on permute les masses externes ;
\item une multiplicité de Moonshine est utilisée comme valeur de masse sans Hamiltonien ;
\item l’opérateur de masse ne commute pas avec la symétrie déclarée préservée ;
\item une hauteur de tour est identifiée à un niveau de corde sans opérateur d’entrelacement ;
\item la condition de niveau ou l’invariance par dualité T échoue ;
\item un écran dimensionnel est déclaré théorie M sans construction de ses champs ;
\item l’existence d’un mode de spin deux est utilisée pour redéfinir la gravité interne antérieure.

\end{itemize}
Ces contrôles maintiennent la connexion positive : il existe un état commun et une classification conjointe ; ce qui est interdit, c’est de remplacer les flèches par une analogie numérique.

\Needspace{7\baselineskip}
\subsection{Construction mathématique des classes matérielles à partir des parcours, des signatures et des représentations}
La double réalisation, la naturalité projective, la chaîne \(A_2^{12}\)--Leech--\(V^\natural\)--Monstre, la trace \(t_3\), les identités \(D_{\rm APP}=54\) et \(196884=54+10\cdot3^9\), ainsi que les écrans dimensionnels appartiennent à l’architecture mathématique APP--TRIT--TPK développée précédemment. La formulation de \(\mathfrak J_{\rm mass,exc}\), la séparation des quatre multiplicités et les contrôles précédents explicitent leur fonction au sein d’une théorie autonome des masses. Les formules conformes et de corde sont des réalisations dimensionnelles ultérieures, avec leurs paramètres déclarés ; elles ne sont pas employées pour engendrer rétroactivement APP, TRIT, TPK, le caractère ou les tours.
